\documentclass[10pt]{article}
\usepackage[T1]{fontenc}
\usepackage[utf8]{inputenc}
\usepackage[french]{babel}
\usepackage{amsmath,amssymb,mathrsfs,mathpazo}
\usepackage{geometry}
\geometry{a4paper,margin=2cm}
\usepackage{graphicx,array,tabularx,tabularray,colortbl,multirow,multicol,booktabs,enumitem,pifont,ulem,fancybox,diagbox,cancel}
\usepackage{pstricks,pstricks-add,pst-plot,pst-tree,pst-node,pst-text,pst-func,pst-math,pst-bspline,pst-3dplot,pst-eucl,pst-solides3d}
\let\pscalculate\undefined
\usepackage{pst-bezier}
\usepackage{tikz,pgfplots,tkz-tab}
\usetikzlibrary{arrows,arrows.meta,patterns,calc,positioning,trees,decorations.markings,shapes.geometric,shapes.misc,shapes.arrows,intersections,decorations.pathmorphing,decorations.pathreplacing,backgrounds,shadows,babel}
\pgfplotsset{compat=1.18}
\usepgfplotslibrary{fillbetween}
\usepackage[np]{numprint}
\usepackage{listings,makecell,pas-tableur,twoopt}
\let\barre\undefined
\usepackage[pstricks]{tablvar}
\definecolor{darkgreen}{RGB}{0,100,0}
\providecommand{\red}{\color{red}}
\providecommand{\blue}{\color{blue}}
\providecommand{\green}{\color{green}}
\providecommand{\black}{\color{black}}
\providecommand{\white}{\color{white}}
\providecommand{\gray}{\color{gray}}
\providecommand{\cyan}{\color{cyan}}
\providecommand{\magenta}{\color{magenta}}
\providecommand{\yellow}{\color{yellow}}
\providecommand{\violet}{\color{violet}}
\providecommand{\orange}{\color{orange}}
\providecommand{\euro}{\texteuro}
\providecommand{\EUR}{\texteuro}
\providecommand{\ieme}{\textsuperscript{e}}
\providecommand{\degres}{\ensuremath{^\circ}}
\providecommand{\degree}{\ensuremath{^\circ}}
\providecommand{\vv}[1]{\overrightarrow{#1}}
\providecommand{\Coord}[1]{#1}
\providecommand{\decosix}{}
\providecommand{\starredbullet}{\ensuremath{\bullet}}
\pagestyle{empty}
\setlength{\parindent}{0pt}
\usepackage{adjustbox}\DeclareUnicodeCharacter{2236}{:}\DeclareUnicodeCharacter{2190}{\ensuremath{\leftarrow}}
\providecommand{\R}{\mathbb{R}}
\providecommand{\N}{\mathbb{N}}
\providecommand{\Z}{\mathbb{Z}}
\providecommand{\Q}{\mathbb{Q}}
\providecommand{\D}{\mathbb{D}}
\providecommand{\C}{\mathbb{C}}
\providecommand{\e}{{\mathrm{e}}}
\providecommand{\vect}[1]{\overrightarrow{#1}}
\providecommand{\vectt}[1]{\overrightarrow{\mathrm{#1}}}
\providecommand{\vvt}[1]{\overrightarrow{\mathrm{#1}}}
\providecommand{\np}[2][]{{#2~#1}}
\providecommand{\vv}[1]{{\overrightarrow{#1}}}
\renewcommand{\d}{{\mathrm{d}}}
\providecommand{\degres}{{^\circ}}
\providecommand{\degre}{{^\circ}}
\providecommand{\blue}{\color{blue}}
\providecommand{\red}{\color{red}}
\providecommand{\enskip}{\quad}
\providecommand{\llbracket}{\lbrack\!\lbrack}
\providecommand{\rrbracket}{\rbrack\!\rbrack}
\providecommand{\euro}{{\text{€}}}
\providecommand{\no}{{\text{n°}}}
\providecommand{\interval}[2]{{[#1;#2]}}
\providecommand{\intervFO}[2]{{[#1;#2[}}
\providecommand{\dotfill}{\ldots}

\providecommand{\ds}{\displaystyle}
\providecommand{\textding}[1]{#1}
\providecommand{\Oij}{\(\mathrm{O};\vec{\imath},\vec{\jmath}\)}
\providecommand{\Oijk}{\(\mathrm{O};\vec{\imath},\vec{\jmath},\vec{k}\)}
\providecommand{\Ouv}{\(\mathrm{O};\vec{u},\vec{v}\)}
\providecommand{\intervFF}[2]{\ensuremath{[#1;#2]}}
\renewcommand{\euro}{\ifmmode\text{\texteuro}\else\texteuro\fi}
\DeclareUnicodeCharacter{00B0}{\ifmmode^{\circ}\else\textdegree\fi}
\DeclareUnicodeCharacter{20AC}{\ifmmode\text{\texteuro}\else\texteuro\fi}
% Énoncés : Éducation nationale. Transcription et corrigés : APMEP, Denis Vergès et contributeurs. Découpage et classement : CoopMaths, Sébastien Lozano, Éric Elter et contributeurs. Adaptation de la présentation : Maths Explorer.
% CC BY-SA : https://coopmaths.fr/www/about/
% https://coopmaths.fr/alea/static/e3c/2020/tex/e3c_2020_00_sujet27_3.tex
\begin{document}
\noindent\textit{APMEP / CoopMaths. CC BY-SA. Présentation adaptée par Maths Explorer.}\par\medskip
\medskip

Dans cet exercice et si cela est nécessaire, les résultats seront arrondis à 0,1.

\begin{minipage}[]{6.25cm}
Le graphique ci-contre illustre l'évolution du nombre (en milliers) de voitures électriques immatriculées en France entre 2015 et 2018.
\end{minipage}
\begin{minipage}[]{8.8125cm}
\psset{xunit=1.5cm,yunit=0.15cm}
\begin{pspicture}(-0.9,-7)(4.85,40)
\psaxes[axesstyle=none,ticks=none,Ox=2015,Dy=5]{-}(0,0)(0,-1)(4.25,33.5)
\multido{\n=0+5}{8}{\psline[linewidth=0.2pt](0,\n)(5,\n)}
\psframe[fillstyle=solid,fillcolor=lightgray](0.75,0)(1.25,17.3)
\psframe[fillstyle=solid,fillcolor=lightgray](1.75,0)(2.25,21.8)
\psframe[fillstyle=solid,fillcolor=lightgray](2.75,0)(3.25,24.9)
\psframe[fillstyle=solid,fillcolor=lightgray](3.75,0)(4.25,31.1)
\uput[u](1,17.3){$17,3$}
\uput[u](2,21.8){$21,8$}
\uput[u](3,24.9){$24,9$}
\uput[u](4,31.1){$31,1$}
\end{pspicture}
\end{minipage}

\medskip

\begin{enumerate}
\item  On cherche à modéliser l'évolution du nombre (en milliers) de voitures électriques immatriculées en France à compter de l'année 2015 à l'aide d'une suite. On hésite entre deux modèles :
\begin{itemize}
\item Premier modèle : 

on fait l'hypothèse que ce nombre augmente de 21\,\% par an. On définit alors une suite $\left(u_n\right)$ où, selon ce modèle, $u_n$ est le nombre (en milliers) de voitures électriques immatriculées en France l'année $2015+n$ avec $n\in \N$. Ainsi, on a $u_0=17,3$.
\item Second modèle :

on définit la suite $\left(v_n\right)$ par $v_0= 17,3$ et pour tout entier naturel $n$,
 \mbox{$v_{n+1}= 0,7v_n + 10$}. D'après ce modèle et pour tout entier naturel $n$, $v_n$ est le nombre (en milliers) de voitures électriques immatriculées en France l'année $2015+n$.
\end{itemize}

\begin{enumerate}
\item Donner les valeurs des réels $u_1$, $u_2$, $u_3$, $v_1$, $v_2$ et $v_3$.
\item Des deux modèles, lequel apparaît le mieux adapté pour modéliser à l'aide d'une suite l'évolution du nombre de voitures électriques immatriculées en France à compter de l'année 2015 donnée dans le graphique ? Argumenter.
\end{enumerate}
\item Dans ce qui suit, on choisit de modéliser le nombre de voitures immatriculées en France à compter de l'année 2015 à l'aide de la suite $\left(u_n\right)$ définie dans la question \textbf{1}.
	\begin{enumerate}
		\item  Donner la nature de la suite $\left(u_n\right)$ et préciser sa raison.
		\item Pour tout entier naturel $n$, exprimer $u_n$ en fonction de $n$.
		\item On considère l'algorithme en langage Python ci-dessous.
\begin{center}
\begin{lstlisting}[style=python]
u=17.3
n=0
while u<50 :
  u=1.21*u
  n=n+1
\end{lstlisting}
\end{center}

Quelle est la valeur de la variable $n$ à la fin de l'exécution de cet algorithme ? Interpréter ce
résultat dans le contexte de l'exercice.
	\end{enumerate}
\end{enumerate}

\vspace{0,5cm}

\end{document}
